\chapter{Introduction}
\label{chp:introduction}

\section{Background}

MicroRNAs (miRNAs), a class of short non-coding RNA molecules \cite{bartel2004microrna}, have emerged as key players in the intricate landscape of post-transcriptional gene regulation. Their role in fine-tuning gene expression is pivotal, impacting various cellular processes, developmental stages, and pathological conditions \cite{bartel2004microrna,balasubr2013developm}. Aberrant expression of miRNAs is related to the progression of many diseases through miRNA-target interactions (MTIs). Databases such as miRTarBase \cite{huang2022mirtarba} and TarBase \cite{karagkou2018diana} have collected thousands of experimental validated MTIs to illustrate the complexity of miRNA regulations in various diseases. Numerous studies have been conducted to survey the potential medical value including biomarker of miRNAs.

The identification and characterization of miRNA biomarkers have opened up exciting avenues in diagnostics and therapeutics. MiRNA regulation is an essential component of gene expression control. These small RNAs bind to the 3\textquotesingle{} untranslated region (UTR) of target mRNAs, leading to translational repression or mRNA degradation. In this manner, miRNAs serve as fine-tuners of protein expression, contributing to the homeostasis of numerous biological systems. The remarkable versatility of miRNAs, combined with their tissue-specific expression patterns, makes them potent candidates for biomarkers in various diseases \cite{krol2010widespre}.

As the field of miRNA regulation has advanced, miRNA biomarkers have gained attention for their potential in disease diagnosis, prognosis, and therapy response prediction. MiRNA expression profiles have been linked to a plethora of pathologies, including cancer, cardiovascular diseases, neurological disorders, and viral infections \cite{mori2019extracel}. Their relatively stable presence in various body fluids, such as blood and urine, underscores their promise as non-invasive diagnostic tools. Researchers have identified that changes in levels of circulating miRNAs have been associated with a wide range of diseases, including type 2 diabetes (T2D), obesity, cardiovascular disease (CVD), cancer, neurodegenerative disorders, and others \cite{mori2019extracel}. Moreover, miRNA dysregulation is associated with different stages of disease progression, offering insights into the prognosis and potential therapeutic targets \cite{zhou2018mirnas}.

The downstream regulatory function of miRNAs, their ability to post-transcriptionally repress target mRNAs, has been the primary focus of miRNA research for over two decades \cite{georgaki2014microtss}. This sustained attention has yielded a wealth of experimentally validated data, curated into comprehensive databases that serve as essential resources for the research community. In contrast to the well-characterized downstream targets, the upstream transcriptional regulation of miRNAs remains less thoroughly explored. Like protein-coding genes, miRNA genes are transcribed by RNA Polymerase II from promoter regions that contain binding sites for transcription factors (TFs) \cite{li2020primirts}. These TF-miRNA regulatory relationships form the "input layer" of the miRNA regulatory network, determining when, where, and at what level specific miRNAs are expressed. Several databases have been developed to catalog TF-miRNA interactions. TransmiR provides a manually curated collection of TF-miRNA regulatory relationships derived from literature mining and experimental evidence \cite{suzuki2002dbtss}. Additionally, resources such as CircuitsDB \cite{mundade2014role} offer computationally predicted TF-miRNA interactions based on sequence motif analysis and regulatory network inference. However, these resources face significant limitations. The accurate identification of miRNA transcription start sites (TSSs), the precise genomic positions where transcription initiates, remains technically challenging \cite{grigoria2022deeptss}. Many current approaches rely on direct mapping of ChIP-seq peaks to miRNA loci, which suffers from high false-positive rates due to signal noise and lacks the resolution necessary for precise functional site delineation \cite{wingende2008transfac}. This incomplete annotation of upstream regulatory elements represents a critical bottleneck in our understanding of miRNA transcriptional control.

The traditional reductionist approach to miRNA research, focusing on individual miRNAs and their specific targets, has generated invaluable insights into discrete regulatory mechanisms \cite{kozomara2019mirbase}. However, this molecule-centric paradigm fails to capture the emergent properties that arise from the collective behavior of miRNAs within complex regulatory networks. The system biology perspective offers a complementary framework that views miRNAs not as isolated entities but as nodes within interconnected regulatory networks. From this viewpoint, biological function is encoded not only in the expression levels of individual molecules but also in the topology of their interactions \cite{wingende2008transfac}. miRNAs occupy a particularly important position within these networks, functioning as regulatory hubs that integrate upstream signals from transcription factors and disperse regulation to hundreds of downstream targets \cite{zou2024chip}.

This network perspective has several advantages. First, network-level properties such as connectivity, modularity, and centrality can reveal functionally important regulatory elements that might be overlooked in single-molecule analyses. Second, network topology tends to be more conserved and robust than individual molecular abundances, potentially offering more stable biomarkers for clinical applications. Third, network analysis can uncover emergent regulatory mechanisms, such as feed-forward loops and feedback circuits, that coordinate gene expression programs across multiple pathways. Despite these conceptual advantages, the practical implementation of network-based miRNA analysis remains limited. Current research predominantly adopts a node-centric approach, focusing on expression level changes while underutilizing the information encoded in network edges, the regulatory relationships between molecules.

While miRNAs are widely recognized for their regulatory roles, the emergence of oxidative modification on miRNAs has introduced an intriguing layer of complexity. Oxidative modification refers to the addition of oxidative moiety 8-hydroxyguanosine (8OH-G), to miRNA nucleotides \cite{depaolis2021epitrans}. The ROS-induced G-to-U transversion of miRNAs, especially in the seed sequence, leads to the altered mRNA targeting \cite{wang2015oxidativ,seok2020position}. Such miRNA modifications alter the original mRNA targeting efficacy (from G-C to U-A pairing), thus allowing the 8OH-G miRNAs to recognize new mRNA target sites with altered targetomes, and significantly impact mRNA stability, target specificity, and regulatory functions \cite{wang2015oxidativ}. 8OH-G miRNA was firstly identified cardiovascular diseases, in the H9c2 rat heart cell line, the ROS-induced 8OH-G modification in miRNA-184 leads to mismatching with the 3' UTR of Bcl-xL and Bcl-w; this may contribute to increased cell apoptosis and enhanced susceptibility of the mouse heart to ischemia/reperfusion injury \cite{wang2015oxidativ}. Phenylephrine treatment was also identified to increase 8OH-G modification of miRNA-1, predominantly at position 7 of the miRNA-1 seed sequence in H9c2 cells. Overexpression of 8OH-G-miRNA-1 or 7-U-miRNA-1 (in which G at position 7 is substituted with U) alone is sufficient to cause cardiac hypertrophy in mice. On the other hand, inhibition of 8OH-G-miRNA-1 in mouse cardiomyocytes attenuated cardiac hypertrophy \cite{seok2020position}. However, the study of miRNA oxidative modification is still in its infancy, and the full extent of its implications is only beginning to be understood. It is known that oxidative stress, a common hallmark of many diseases, can induce the modification. Consequently, the exploration of miRNA oxidative modification has significant implications for our understanding of the role of oxidative stress in disease pathogenesis.

Recent advancements in sequencing technologies and innovative experimental approaches have paved the way for the systematic investigation of miRNA oxidative modification, by combining immunoprecipitation and small RNA sequencing, the mystery of oxidative miRNA have been enlightened \cite{seok2020position}. The identification of specific modified miRNAs associated with diseases could potentially serve as valuable diagnostic and therapeutic targets. Moreover, understanding the regulatory effects of miRNA oxidative modification will enhance our comprehension of disease mechanisms and potentially offer novel strategies for intervention.

\section{Motivation and Research Gaps}

The first critical gap in current miRNA research concerns the incomplete annotation of upstream transcriptional regulation. Accurate identification of miRNA transcription start sites (TSSs) is fundamental to understanding how miRNA expression is controlled, as TSSs anchor the promoter regions where transcription factors bind to initiate transcription. However, TSS identification for miRNAs presents unique challenges compared to protein-coding genes. The genomic architecture of miRNA genes is highly variable: approximately 50\% of human miRNAs are intragenic, located within introns of host genes and potentially sharing regulatory elements with their hosts, while the remaining intergenic miRNAs possess independent promoters that can be located tens of kilobases upstream from the mature miRNA sequence. This architectural diversity complicates the systematic annotation of miRNA TSSs.

Existing experimental approaches, such as cap analysis of gene expression (CAGE-seq), have been successfully applied to map TSSs for mRNAs but face limitations when applied to miRNAs. The relatively low expression levels of primary miRNA transcripts (pri-miRNAs) compared to mRNAs result in reduced signal-to-noise ratios, compromising detection sensitivity. Computational methods have attempted to predict miRNA TSSs using sequence features and epigenetic marks, but models trained on limited datasets often suffer from high false-positive rates. Consequently, current TF-miRNA interaction databases are constrained by the absence of accurate TSS annotations. Without knowing precisely where transcription begins, the mapping of TF binding sites to functional promoter regions remains imprecise, limiting our ability to construct comprehensive upstream regulatory networks.

The second critical gap concerns a systematic bias in how miRNA networks are analyzed. Despite the conceptual shift toward systems biology, most current research remains fundamentally node-centric, focused on measuring and comparing the expression levels of individual miRNAs across conditions. This node-centric approach has practical origins: expression profiling is technically straightforward and generates readily interpretable data. However, it overlooks a fundamental aspect of network biology: information is encoded not only in nodes (molecules) but also in edges (relationships between molecules). The regulatory interaction between a miRNA and its target, including the strength, context, and functional consequence of that interaction, carries information that is distinct from and complementary to the expression level of either partner.

Edge information operates at two layers. At the first layer, individual edges represent specific miRNA-target interactions. The regulatory coupling between a miRNA and its target --- whether the interaction is active, what regulatory strength it exerts --- can provide clinical information that surpasses expression levels alone. At the second layer, edge patterns represent the overall topology of how a miRNA connects to different targets and pathways. This topological context can determine the functional role of a miRNA in ways that its expression level cannot predict. Current research largely underutilizes both layers of edge information. Prognostic models typically rely on expression signatures, missing the potentially more robust signal encoded in regulatory relationships. Functional interpretation focuses on target identification rather than network architecture, failing to explain context-dependent phenomena such as why the same miRNA can exhibit opposite functions in different disease settings.

The third critical gap concerns the epitranscriptomic dimension of miRNA regulation, specifically, the role of oxidative modifications in dynamically rewiring miRNA regulatory networks. Under conditions of oxidative stress, which are prevalent in pathological states including cancer, reactive oxygen species (ROS) can chemically modify guanine bases in RNA to form 8-oxo-guanine (8-oxoG). When this modification occurs within the seed region of a miRNA --- the 7-nucleotide sequence critical for target recognition --- it alters the base-pairing preference from canonical G:C pairing to non-canonical 8-oxoG:A pairing. This single nucleotide change can fundamentally redirect the miRNA to a new set of targets, effectively rewiring the regulatory network without any change in miRNA expression level.

Despite the profound implications of this mechanism, systematic investigation of miRNA oxidation has been severely limited by technical barriers. The current gold-standard approach, immunoprecipitation with anti-8-oxoG antibodies followed by sequencing (8-oxoG-IP-seq) is technically demanding, expensive, and requires substantial input material. These constraints have resulted in fewer than 10 published studies systematically examining miRNA oxidation, leaving the prevalence, specificity, and functional consequences of this modification largely unexplored. Furthermore, the absence of standardized computational pipelines for analyzing oxidation data has hindered reproducibility and comparability across studies. The lack of publicly available oxidation profiles across diverse biological contexts represents a significant gap in our understanding of epitranscriptomic regulation, limiting the translation of this mechanism into clinical applications.

\section{Objectives of the Study}

\begin{enumerate}
\item \textbf{Decoding transcriptional control and establishing upstream regulatory data foundation.} The first aim of this thesis is to develop a deep learning platform and comprehensive regulatory atlas for miRNA transcription start sites. Integrate multi-modal sequencing data including CAGE-seq, TSS-seq, ChIP-seq, and DNase-seq to train a convolutional neural network model while systematically annotating transcription factor binding sites in predicted promoter regions, enabling tissue-specific exploration through an interactive web database with genome browser visualization and conservation analysis for targeted experimental design like CRISPR modulation.

\item \textbf{Elucidating network architecture and demonstrating edge-centric value.} The second aim of this thesis is to demonstrate the compensatory effect of edge-centric miRNA network analysis. Moreover, developing a framework quantifying miRNA-target interaction strength for clinical modeling in triple-negative breast cancer while reconstructing multi-layer networks for hub miRNAs incorporating upstream and downstream elements, enabling demonstration of complementary insights from individual edges and patterns that surpass node-centric approaches for dynamic network understanding.

\item \textbf{Unveiling oxidative rewiring effect and extending the network dynamics.} The third aim of this thesis is to develop a computational pipeline and database for detecting oxidative miRNA modifications. Construct an oxidized reference library with oxidation-aware alignment strategies benchmarked against experimental data and spike-in controls while quantifying rewiring through a Rewiring Activity Score for target switching analysis, enabling identification of prognostic biomarkers and mechanistic links to disease phenotypes across species for oxidative modification research.
\end{enumerate}

Collectively, this thesis is intended to create a systematic framework progressing from transcriptional regulation to epitranscriptomic modification. By integrating deep learning, network biology, and oxidative modification into a unified analytical pipeline, enabling comprehensive dissection of miRNA-mediated gene regulation across multiple dimensions and advancing both methodological innovation and mechanistic understanding.
